% =============================================================================
%  CHAPTER 1 — INTRODUCTION
%
%  >>> MOTIVATION + REPORT STRUCTURE are an AI-ASSISTED FIRST DRAFT (from your
%      docs/PROJECT_CONTEXT.md + RELATED_WORK.md). Rewrite in your own voice,
%      verify citations, and disclose in the Declarations page.
%  >>> OBJECTIVES, RESEARCH-QUESTION FRAMING, and CONTRIBUTIONS are left for YOU
%      — they are your claims and the heart of the Framing mark (15%). I'll edit
%      what you write, but these should not be ghost-drafted.
%  RUBRIC: Framing of Research Problem (15%). BUDGET: ~4-6 pages.
% =============================================================================
\chapter{Introduction}
\label{ch:intro}

\section{Motivation}
\label{sec:intro-motivation}

Software spends most of its life being maintained rather than written, and a
large share of that maintenance is \emph{refactoring}: behaviour-preserving
restructuring that keeps a codebase readable and able to evolve. The arrival of
capable large language models (LLMs) inside the IDE has changed how developers
approach such work --- controlled studies report that AI pair-programmers let
developers complete tasks substantially faster \citep{copilotproductivity2023}
--- and LLMs are now demonstrably able to \emph{identify} useful refactorings
and describe how to carry them out \citep{cordeiro2024refactoring,
potentialllm2024}.

The difficulty lies not in deciding \emph{what} to refactor but in
\emph{applying} the change correctly. When an LLM performs a refactoring by
generating modified source text, it must itself locate and update every affected
site. In practice it routinely misses usages in other files
(\emph{scope blindness}), breaks import statements, introduces compilation
errors, or produces code that compiles but is semantically wrong
\citep{llmrefactorlimits2025}. These are not occasional lapses: on a
repository-level benchmark of 1{,}099 real, behaviour-preserving Java
refactorings, the strongest evaluated LLM reaches a pass rate of only 41.6\%
\citep{swerefactor2025}. The model is reliable at \emph{what} and unreliable at
\emph{how} --- precisely because ``how'' has been framed as the problem of
writing correct code text.

Yet that problem is already solved, just not by the model. Every modern IDE
contains a refactoring engine that performs renames, moves, signature changes,
and extractions over a typed program model, updating all references, imports,
and occurrences atomically and correctly by construction. The weakness of
current LLM agents is that they do not use it for writes: agent toolkits expose a
general-purpose textual \texttt{edit} action \citep{sweagent2024}, and even
systems that exploit the abstract syntax tree do so only to \emph{read} context
before emitting a text patch \citep{autocoderover2024}. Where the IDE engine
\emph{is} used for execution, as in EM-Assist \citep{emassist2024}, it is wired
to a single operation inside the IDE, with no interface an external agent could
call.

This project argues that an LLM agent should be \emph{decoupled} from the act of
writing code: the model should decide \emph{what} to refactor, and the IDE's
refactoring engine should decide and execute \emph{how}. Concretely, it presents
an IntelliJ~IDEA plugin that exposes the IDE's structured refactoring operations
--- rename, move, change-signature, safe-delete, inline, extract, and symbol
creation --- through a localhost tool API that any function-calling LLM can
drive. Because every mutation is performed by IntelliJ's PSI engine rather than
by text substitution, the result is correct with respect to the IDE's semantic
model independently of the quality of the model's textual output.
This report defends a single claim: that the IDE's structured operations are a
better action space for an LLM refactoring agent than source text --- not because
a capable model cannot edit text correctly, but because routing every edit
through IntelliJ's PSI engine makes project-wide, behaviour-preserving refactoring
of real-world Java codebases correct \emph{by construction} and more efficient,
and uniquely supports the project-wide structural queries (such as exhaustive
reference finding) that text editing cannot perform. The advantage is largest
exactly where text editing is least reliable: under weaker models, and for
refactorings whose correctness depends on a project-wide reference or type graph.

\section{Objectives}
\label{sec:intro-objectives}
% YOUR SECTION. Turn your research questions (docs/EVALUATION.md) into explicit,
% testable objectives. List only what the evaluation actually measures; move the
% developer study (RQ4) to Future Work if you do not run it.
This project pursues four objectives, each phrased so that it is tested directly
by the evaluation in Chapter~\ref{ch:eval}.

\begin{enumerate}
  \item \textbf{Expose the IDE's refactoring engine as an agent action space.}
  Design and implement an IntelliJ~IDEA plugin that makes the platform's
  PSI-backed operations --- rename, move, change-signature, safe-delete, inline,
  extract, and symbol creation --- callable by an external function-calling LLM
  through a localhost tool API.
  \emph{(Substantiated in Chapters~\ref{ch:design}--\ref{ch:impl}.)}

  \item \textbf{Demonstrate behaviour-preserving correctness on real-world code.}
  Show that refactorings executed through this action space are applied
  correctly on two real-world Java projects (spring-petclinic and Apache Commons
  Lang), measured by per-task pass rate under compilation and symbol-survival
  checks. \emph{(Chapter~\ref{ch:eval}.)}

  \item \textbf{Establish where structured execution is \emph{necessary}.}
  Isolate and measure the refactoring classes on which text-level editing fails
  --- cross-file moves requiring import repair, overload-safe and cross-class
  renames, and project-wide usage discovery --- to show where the structured
  approach is not merely tidier but required for correctness.
  \emph{(Chapter~\ref{ch:eval}, root-cause analysis.)}

  \item \textbf{Compare against a text-editing baseline.}
  Evaluate the structured approach against both a scripted text editor and an LLM
  text-edit agent on the same tasks, comparing correctness and the minimality of
  the resulting diffs. \emph{(Chapter~\ref{ch:eval}.)}
\end{enumerate}

\section{Research-Question Framing}
\label{sec:intro-rqs}
% YOUR SECTION. Frame the work as AI-systems research, not just a plugin: e.g.
% "how should an LLM agent's action space be designed for safe, reliable,
% project-wide code transformation?" Cite the ACI abstraction \citep{sweagent2024}
% and the constrained-action-space taxonomy \citep{demystifying2025}.
We frame this work not as a plugin but as a question about \emph{agent design}.
Under the agent--computer-interface (ACI) abstraction \citep{sweagent2024}, an
LLM agent's competence is governed less by raw model capability than by the
\emph{action space} it is given --- the set of operations through which it
observes and changes the world; recent analysis of code agents makes the same
point, treating the design of a constrained, well-typed action space as a
primary determinant of reliability \citep{demystifying2025}. Conventional coding
agents take \emph{source text} as their action space. This work asks what
changes when the action space is instead the IDE's structured refactoring
engine:

\begin{description}
  \item[RQ1 (Correctness).] Does constraining an agent's action space to the
  IDE's structured, AST-aware refactoring operations --- rather than free-form
  text edits --- improve the correctness of behaviour-preserving, project-wide
  refactorings on real-world codebases? \emph{(Objectives~2, 4.)}

  \item[RQ2 (Necessity).] For which classes of refactoring is a structured
  action space not merely beneficial but \emph{necessary} --- i.e.\ where does
  text-level editing fail for reasons intrinsic to its action representation
  rather than to model capability? \emph{(Objective~3.)}

  \item[RQ3 (Cost).] What does this safety cost? Does routing every edit through
  the IDE engine change the minimality of the resulting diffs and the agent's
  interaction cost (tool calls and turns) relative to a text-editing agent?
  \emph{(Objective~4.)}
\end{description}

\section{Contributions}
\label{sec:intro-contributions}
% YOUR SECTION. Adapt the novelty claims in docs/PROJECT_CONTEXT.md, but HEDGE
% absolute "first/only" claims to "to the best of our knowledge" and acknowledge
% contemporaneous work (e.g. the JetBrains MCP server). Each contribution should
% map to a chapter/section.
The contributions of this report are:

\begin{enumerate}
  \item \textbf{A structured-refactoring tool API for LLM agents.} An
  IntelliJ~IDEA plugin that exposes the platform's PSI-backed refactoring
  operations through a model-agnostic, function-calling HTTP interface, letting
  an external LLM \emph{execute} --- not merely propose --- project-wide
  refactorings. To the best of our knowledge this is the first general-purpose
  interface of its kind for arbitrary function-calling agents; it is
  contemporaneous with, and complementary to, the JetBrains MCP server.
  \emph{(Chapters~\ref{ch:design}--\ref{ch:impl}.)}

  \item \textbf{A headless-execution technique for IDE refactorings.} A method
  for making \texttt{BaseRefactoringProcessor}-based refactorings fully
  non-interactive --- suppressing the preview, conflict, and error dialogs that
  otherwise block an unattended IDE --- so that an agent can drive them without
  a human at the keyboard. \emph{(Chapter~\ref{ch:impl}.)}

  \item \textbf{An evaluation methodology and benchmark.} A reproducible harness
  that compares structured execution against both a scripted text editor and an
  LLM text-edit agent on real-world Java projects, scoring correctness
  (compilation and symbol survival) and diff minimality.
  \emph{(Chapter~\ref{ch:eval}.)}

  \item \textbf{An empirical characterisation of where the action space matters.}
  Evidence that text-level editing fails on reference-graph- and type-dependent
  refactorings --- cross-file moves with import repair, overload-safe and cross-class
  renames, and project-wide usage discovery --- for scripted and weaker-model drivers;
  that a \emph{capable} agent reaches parity on those refactorings; and that the
  structured engine's durable advantages are therefore correctness by construction, the
  structural queries text editing cannot express, and robustness to the driver's
  capability. \emph{(Chapter~\ref{ch:eval}.)}
\end{enumerate}

\section{Report Structure}
\label{sec:intro-structure}

The remainder of this report is organised as follows.
Chapter~\ref{ch:background} reviews the relevant literature on LLM-driven
refactoring, program-transformation agents, AST-safe refactoring, and tool use,
and identifies the gap this work addresses.
Chapter~\ref{ch:design} presents the design: the decoupling of model reasoning
from IDE execution, the tool API, and the rationale behind the key design
decisions.
Chapter~\ref{ch:impl} describes the implementation, including how each operation
delegates to IntelliJ's refactoring engine and the principal engineering
challenges overcome.
Chapter~\ref{ch:eval} evaluates the system against a text-editing baseline on
two benchmark projects, reporting correctness, efficiency, and a root-cause
analysis of the baseline's failures.
Chapter~\ref{ch:conclusion} concludes with a reflection on the results, their
broader impact, and directions for future work.
